\subsection{Pseudo-image measures}

DEFINITION 1. --- Let T and B be two locally compact spaces, \(\mu\) a positive measure on T, and p a \(\mu\) -measurable mapping of T into B. A positive measure \(\nu\) on B is said to be a pseudo-image measure of \(\mu\) under p if it satisfies the following condition: for a subset N of B to be locally \(\nu\) -negligible, it is necessary and sufficient that \(\overline{p}^{-1}(N)\) be locally \(\mu\) -negligible.

Examples. --- 1) If \(p\) is \(\mu\) -proper and \(\nu = p(\mu)\) , then \(\nu\) is a pseudo-image measure of \(\mu\) under \(p\) (Ch. V, §6, No.~2, Cor. 2 of Prop. 2).

\begin{enumerate}
\def\labelenumi{\arabic{enumi})}
\setcounter{enumi}{1}
\tightlist
\item
  Let \(B'\) be a locally compact space, \(\nu'\) a positive measure on \(B'\) ; take for T the space \(B \times B'\) and for \(\mu\) the measure \(\nu \otimes \nu'\) ; if p is the projection of T onto B, then \(\nu\) is a pseudo-image of \(\mu\) under p (Ch. V, §8, No.~2, Prop. 4 and No.~3, Cor. 1 of Prop. 7).
\end{enumerate}

Note that if \(\nu\) is a pseudo-image measure of \(\mu\) under \(p\) , then \(\nu\) is carried by \(p(T)\) .

If \(\nu\) is a pseudo-image of \(\mu\) under p, the set of measures that are pseudo-images of \(\mu\) under p is the class of positive measures equivalent to \(\nu\) , and every positive measure equivalent to \(\mu\) admits the same pseudo-image measures under p.~The class of \(\nu\) is said to be the pseudo-image class of that of \(\mu\) under p.

PROPOSITION 1. --- Let T be a locally compact space countable at infinity, \(\mu\) a positive measure on T, and p a \(\mu\) -measurable mapping of T into a locally compact space B. Then there exists a pseudo-image measure of \(\mu\) under p.

For, there exists a bounded measure \(\mu'\) on T equivalent to \(\mu\) (Ch. V, §5, No.~6, Prop. 11); \(p\) is then \(\mu'\) -proper.

